% =====================================================================
\section*{Executive summary}
\addcontentsline{toc}{section}{Executive summary}

\textbf{Question.} Where will a surface drifter in the Gulf Stream be 24 hours from now, and does knowing how it moved over the previous 24 hours improve that prediction beyond simply extrapolating its present velocity? We also ask under which recent motion conditions---speed, acceleration and turning---24-hour forecasts are most difficult, and where trajectory history helps most.

\textbf{Data and design.} Using the NOAA Global Drifter Program (GDP) hourly product \citep{elipot2016,gdpdata2022}, we built \num{626679} forecast cases from \num{517} drifters with forecast origins in the box \ang{80}--\ang{60}W, \ang{30}--\ang{45}N during July--September 2007--2022. Each case uses 25 hourly observations ($t-24$\,h to $t$) to predict the observed position at $t+24$\,h. Two physical baselines (Constant Velocity and empirical Mean-flow Advection) were compared with five learned model families (Ridge, Elastic Net, Random Forest, Extra Trees and Gradient Boosting), each fitted with \emph{current-state} features only and with additional \emph{history} features. Learned models predict a correction to the Constant Velocity displacement, which is converted to a geographic position by spherical geodesic formulas. Drifters were split once into disjoint training (361), validation (78) and test (78) sets; hyperparameters were chosen on validation data and the frozen models were evaluated once on the \num{82546} held-out test forecasts.

\textbf{Main findings (held-out test set).}
\begin{itemize}
  \item History-enhanced Gradient Boosting (\GBH) had the lowest error: median \SI{9.23}{km} and 90th percentile (P90) \SI{28.55}{km}, against \SI{15.89}{km} and \SI{40.90}{km} for Constant Velocity (reductions of 41.9\% and 30.2\%). Random Forest and Extra Trees with history performed almost identically (medians \SI{9.25}{km} and \SI{9.33}{km}); we do not claim a ranking among the three.
  \item Adding 24-hour history reduced median error in every model family, by 17--18\% for the linear models and 29\% for the tree ensembles. For Gradient Boosting, the mean paired error reduction was \SI{3.65}{km} (95\% drifter-cluster bootstrap interval \SIrange{3.14}{4.14}{km}).
  \item The Mean-flow Advection baseline, built from a training-data climatology, was the least accurate method overall (median \SI{24.28}{km}).
  \item Errors grow with recent speed and acceleration but \emph{fall} with turning rate, largely because fast drifters tend to move straight and slow drifters turn more. The benefit of history was positive in all nine speed--turning regimes and largest for fast, strongly turning drifters (\SI{7.03}{km}, 95\% interval \SIrange{5.38}{9.17}{km}).
  \item The advantage of \GBH{} held in every one of the 16 years. Sea-surface temperature was associated with forecast difficulty, but adding it as a predictor changed validation error by less than 0.2\%.
\end{itemize}

\textbf{Why it matters.} Shrinking the typical 24-hour miss from about \SI{16}{km} to \SI{9}{km} is the kind of gain that search-and-rescue, pollution-tracking and ocean-monitoring users of drift estimates care about. Equally useful is knowing \emph{in advance} which forecasts are unreliable: recent speed, which is available at forecast time, is strongly associated with larger errors, whereas track curvature is not a warning sign. These applications were not tested here, and other floating objects drift differently from drogued buoys.

\textbf{Recommendation.} For the industry partner and other users of the GDP hourly product, a history-enhanced tree-ensemble correction to Constant Velocity is the recommended benchmark for \emph{retrospective} evaluation of 24-hour drifter predictability, and is clearly preferable to either physical baseline. Gradient Boosting is the practical choice because it was by far the fastest of the three to train. GDP hourly velocities are smoothed using fixes on both sides of each time, and almost half of the forecast origins come from drifters that had lost their drogue. The results are therefore \emph{not} yet evidence of operational real-time forecasting skill. Recommended next steps are evaluation with real-time inputs, drogue-stratified and chronological hold-outs, multiple lead times and, for applications, validation against the drift of the relevant objects.
